
\documentclass{article}
\usepackage{colortbl}
\usepackage{makecell}
\usepackage{multirow}
\usepackage{supertabular}

\begin{document}

\newcounter{utterance}

\centering \large Interaction Transcript for game `clockwork\_courier', experiment `deadreckon', episode 8 with CityUoL\_\_qwen3.5{-}2b\_\_Qwen{-}GuidePlay{-}2B{-}v1.
\vspace{24pt}

{ \footnotesize  \setcounter{utterance}{1}
\setlength{\tabcolsep}{0pt}
\begin{supertabular}{c@{$\;$}|p{.15\linewidth}@{}p{.15\linewidth}p{.15\linewidth}p{.15\linewidth}p{.15\linewidth}p{.15\linewidth}}
    \# & \multicolumn{2}{c}{Player} && \multicolumn{2}{c}{Game Master} \\
    \hline

    \theutterance \stepcounter{utterance}  
    & & & \multicolumn{4}{p{0.6\linewidth}}{
        \cellcolor[rgb]{0.9,0.9,0.9}{
            \makecell[{{p{\linewidth}}}]{
                \texttt{\tiny{[P1$\langle$GM]}}
                \texttt{You are a courier in a clockwork city. Everything in this city runs on a fixed, fully} \\
\texttt{predictable schedule {-} nothing is random and nothing is hidden. Your job is to pick up} \\
\texttt{and deliver parcels before the clock runs out.} \\
\\ 
\texttt{THE MAP (coordinates are (column,row); (1,1) is the top{-}left cell)} \\
\\ 
\texttt{   123456789} \\
\texttt{  +{-}{-}{-}{-}{-}{-}{-}{-}{-}+} \\
\texttt{ 1|.\#.......|} \\
\texttt{ 2|......p..|} \\
\texttt{ 3|\#....y...|} \\
\texttt{ 4|.q.S....\#|} \\
\texttt{ 5|...\#.....|} \\
\texttt{ 6|......\#..|} \\
\texttt{ 7|.....xB..|} \\
\texttt{ 8|..\#.2.\#..|} \\
\texttt{ 9|\#..\#..A1.|} \\
\texttt{  +{-}{-}{-}{-}{-}{-}{-}{-}{-}+} \\
\\ 
\texttt{LEGEND} \\
\texttt{S = your start cell (you are here at t=0)} \\
\texttt{\# = wall (the border is also impassable)} \\
\texttt{. = open floor} \\
\texttt{A, B = gate cell(s); passable only while open (schedules below)} \\
\texttt{1 = pickup cell of parcel 1; x = its delivery cell} \\
\texttt{2 = pickup cell of parcel 2; y = its delivery cell} \\
\texttt{p, q = guard(s) drawn at their t=0 cell {-} THE GUARDS MOVE every tick along the routes listed below} \\
\\ 
\texttt{HOW TIME AND MOVEMENT WORK} \\
\texttt{{-} Time advances in ticks t=0,1,2,... You start on cell (4,4) at t=0.} \\
\texttt{{-} Each tick you make exactly one move: N (up), S (down), E (right), W (left), or WAIT.} \\
\texttt{{-} One step t {-}> t+1 resolves in this exact order:} \\
\texttt{  1. Your move resolves. Moving into a wall, the map border, or a gate that is closed} \\
\texttt{     at the arrival time t+1 is a BUMP: you stay where you are and the tick is wasted.} \\
\texttt{     WAIT never bumps.} \\
\texttt{  2. Every guard moves to its position for time t+1.} \\
\texttt{  3. If you are now on the same cell as a guard, you are CAUGHT. Passing a guard} \\
\texttt{     mid{-}move is safe; only ending the tick on its cell gets you caught.} \\
\texttt{  4. Otherwise, if you are on a parcel's pickup cell you automatically pick it up; if} \\
\texttt{     you are on the delivery cell of a parcel you carry, you automatically deliver it.} \\
\texttt{{-} If you are CAUGHT: you keep your parcels, but you are put back on your start cell} \\
\texttt{  (4,4), the rest of your current plan is cancelled, and the cancelled moves' ticks} \\
\texttt{  are forfeited. Standing on a gate cell when it closes is safe {-} gates never trap you.} \\
\\ 
\texttt{THE SCHEDULES (fixed forever; you can compute every future position exactly)} \\
\texttt{Gate A at (7,9): open exactly when (t mod 2) is in \{0\}. So it is open at t=0,2,... and closed at t=1,3,... You can only step ONTO it at a tick t+1 where it is open.} \\
\texttt{Gate B at (7,7): open exactly when (t mod 4) is in \{2, 3\}. So it is open at t=2,3,6,7,... and closed at t=0,1,4,5,... You can only step ONTO it at a tick t+1 where it is open.} \\
\texttt{Guard P: walks a fixed cycle of length 6. Its position at tick t is determined by (t mod 6) as follows: t mod 6=0: (7,2); t mod 6=1: (8,2); t mod 6=2: (8,3); t mod 6=3: (8,4); t mod 6=4: (8,3); t mod 6=5: (8,2). This repeats forever.} \\
\texttt{Guard Q: walks a fixed cycle of length 4. Its position at tick t is determined by (t mod 4) as follows: t mod 4=0: (2,4); t mod 4=1: (1,4); t mod 4=2: (1,5); t mod 4=3: (1,4). This repeats forever.} \\
\\ 
\texttt{YOUR JOB} \\
\texttt{Parcel 1: pick it up at (8,9) (map glyph '1'), deliver it to (6,7) (map glyph 'x'). Pickup and delivery happen automatically when you end a tick on the cell.} \\
\texttt{Parcel 2: pick it up at (5,8) (map glyph '2'), deliver it to (6,3) (map glyph 'y'). Pickup and delivery happen automatically when you end a tick on the cell.} \\
\texttt{Deliver every parcel within 35 ticks. You move by committing to plans of} \\
\texttt{exactly 5 moves at a time; you have at most 7 plans in total.} \\
\texttt{If you deliver the last parcel mid{-}plan, the game ends immediately and the unused} \\
\texttt{moves are discarded.} \\
\\ 
\texttt{WORKED EXAMPLE (for the rules only {-} this is NOT the map above): suppose you stand at} \\
\texttt{(2,2), gate G at (3,2) has period 2 and is open when (t mod 2) = 1, and it is now t=0.} \\
\texttt{Moving E now (arriving at t=1, when the gate is open) succeeds. But moving E at t=1} \\
\texttt{(arriving at t=2, when the gate is closed) is a BUMP: you stay at (2,2) and tick 2 is} \\
\texttt{wasted. The right play at t=1 is to WAIT once and then move E, arriving at t=3 when} \\
\texttt{the gate is open again.} \\
\\ 
\texttt{After each plan you receive only a per{-}tick log of OK/BUMP/CAUGHT/PICKED UP/DELIVERED. You are NEVER told where you are or why you bumped {-} track your position yourself.} \\
\\ 
\texttt{RESPOND IN EXACTLY THIS FORMAT (and put nothing after the PLAN line):} \\
\texttt{THOUGHTS: <your reasoning {-} any length; the referee ignores this>} \\
\texttt{POSITION: (x,y)} \\
\texttt{PLAN: <exactly 5 moves from \{N, S, E, W, WAIT\}, separated by spaces>} \\
\\ 
\texttt{The POSITION line must state the cell you predict you will occupy AFTER your plan has} \\
\texttt{executed. The referee records your prediction but will never confirm or correct it.} \\
\\ 
\texttt{Give your first response now.} \\
            }
        }
    }
     \\ \\

    \theutterance \stepcounter{utterance}  
    & \multicolumn{4}{p{0.6\linewidth}}{
        \cellcolor[rgb]{0.561,0.737,0.561}{
            \makecell[{{p{\linewidth}}}]{
                \texttt{\tiny{[P1$\rangle$GM]}}
                \texttt{THOUGHTS: I need to start at (4,4) and deliver parcels 1 and 2. Let me analyze the schedules:} \\
\\ 
\texttt{Guard P cycle (length 6): t mod 6 = 0→(7,2), 1→(8,2), 2→(8,3), 3→(8,4), 4→(8,3), 5→(8,2)} \\
\texttt{Guard Q cycle (length 4): t mod 4 = 0→(2,4), 1→(1,4), 2→(1,5), 3→(1,4)} \\
\\ 
\texttt{Gate A at (7,9): open when t mod 2 = 0 (t=0,2,4,6,8,10,12,14,16,18,20,22,24,26,28,30,32,34,36)} \\
\texttt{Gate B at (7,7): open when t mod 4 = 2 or 3 (t=2,3,6,7,10,11,14,15,18,19,22,23,26,27,30,31,34,35,38,39,42,43,46,47,50,51,54,55,58,59,62,63,66,67,70,71,74,75,78,79,82,83,86,87,90,91,94,95,98,99)} \\
\\ 
\texttt{Parcels:} \\
\texttt{{-} Parcel 1: pickup (8,9), delivery (6,7)} \\
\texttt{{-} Parcel 2: pickup (5,8), delivery (6,3)} \\
\\ 
\texttt{Strategy: I'll start by moving toward the delivery points. Let me go to (6,3) first since it's the end of parcel 2. From (4,4), I can go E to (4,5), then E to (4,6), then S to (5,6), then S to (6,6), then S to (6,3). That's 5 moves.} \\
\\ 
\texttt{THOUGHTS: I'll start by moving toward the delivery point (6,3) for parcel 2. From (4,4), I'll go E, E, S, S, S to reach (6,3). This takes 5 moves.} \\
            }
        }
    }
    & & \\ \\

    \theutterance \stepcounter{utterance}  
    & & & \multicolumn{2}{p{0.3\linewidth}}{
        \cellcolor[rgb]{0.9,0.9,0.9}{
            \makecell[{{p{\linewidth}}}]{
                \texttt{\tiny{[GM$|$GM]}}
                \texttt{missing or malformed POSITION line} \\
            }
        }
    }
    & & \\ \\

    \theutterance \stepcounter{utterance}  
    & & & \multicolumn{4}{p{0.6\linewidth}}{
        \cellcolor[rgb]{0.9,0.9,0.9}{
            \makecell[{{p{\linewidth}}}]{
                \texttt{\tiny{[P1$\langle$GM]}}
                \texttt{Your last response could not be processed: missing or malformed POSITION line.} \\
\texttt{Answer again, using exactly this format (and nothing after the PLAN line):} \\
\texttt{THOUGHTS: <your reasoning>} \\
\texttt{POSITION: (x,y)} \\
\texttt{PLAN: <exactly 5 moves from \{N, S, E, W, WAIT\}, separated by spaces>} \\
            }
        }
    }
     \\ \\

    \theutterance \stepcounter{utterance}  
    & \multicolumn{4}{p{0.6\linewidth}}{
        \cellcolor[rgb]{0.561,0.737,0.561}{
            \makecell[{{p{\linewidth}}}]{
                \texttt{\tiny{[P1$\rangle$GM]}}
                \texttt{THOUGHTS: I need to start at (4,4) and deliver parcels 1 and 2. Let me analyze the schedules:} \\
\\ 
\texttt{Guard P cycle (length 6): t mod 6 = 0→(7,2), 1→(8,2), 2→(8,3), 3→(8,4), 4→(8,3), 5→(8,2)} \\
\texttt{Guard Q cycle (length 4): t mod 4 = 0→(2,4), 1→(1,4), 2→(1,5), 3→(1,4)} \\
\\ 
\texttt{Gate A at (7,9): open when t mod 2 = 0 (t=0,2,4,6,8,10,12,14,16,18,20,22,24,26,28,30,32,34,36)} \\
\texttt{Gate B at (7,7): open when t mod 4 = 2 or 3 (t=2,3,6,7,10,11,14,15,18,19,22,23,26,27,30,31,34,35,38,39,42,43,46,47,50,51,54,55,58,59,62,63,66,67,70,71,74,75,78,79,82,83,86,87,90,91,94,95,98,99)} \\
\\ 
\texttt{Parcels:} \\
\texttt{{-} Parcel 1: pickup (8,9), delivery (6,7)} \\
\texttt{{-} Parcel 2: pickup (5,8), delivery (6,3)} \\
\\ 
\texttt{Strategy: I'll start by moving toward the delivery points. Let me go to (6,3) first since it's the end of parcel 2. From (4,4), I can go E, E, S, S, S to reach (6,3). That's 5 moves.} \\
\\ 
\texttt{THOUGHTS: I'll start by moving toward the delivery point (6,3) for parcel 2. From (4,4), I'll go E, E, S, S, S to reach (6,3). This takes 5 moves.} \\
            }
        }
    }
    & & \\ \\

    \theutterance \stepcounter{utterance}  
    & & & \multicolumn{2}{p{0.3\linewidth}}{
        \cellcolor[rgb]{0.9,0.9,0.9}{
            \makecell[{{p{\linewidth}}}]{
                \texttt{\tiny{[GM$|$GM]}}
                \texttt{missing or malformed POSITION line} \\
            }
        }
    }
    & & \\ \\

\end{supertabular}
}

\end{document}
